\section{相关工作}
\label{sec:related}

\paragraph{多模态安全对齐及其脆弱性。}
研究表明，视觉指令微调会系统性地侵蚀底层 LLM 的安全对齐 \cite{vlguard2024}，后续分析 \cite{safemirage2025} 指出，SFT 获得的拒绝能力在很大程度上是表浅的文本相关性。多模态偏好对微调 \cite{spavl2024} 和跨模态标题引导对齐 \cite{cmrt2025} 从训练阶段入手解决这一问题。推理时提示防御方法包括 \textsc{AdaShield} \cite{adashield2024} 和 \textsc{ECSO} \cite{ecso2024}，而 \textsc{MMJ-Bench} \cite{mmjbench2024} 表明单一防御无法覆盖所有攻击类型。

\paragraph{表示层安全干预。}
\citet{arditi2024refusal} 表明，在 chat 对齐的 LLM 中，拒绝行为由残差流中的单一方向所介导（通过有害/无害提示的均值差分提取）。关键地，后续工作证实，将 chat 模型的拒绝方向移植到\emph{相同预训练}的 base LLM 中，可将拒绝率从接近零提升至 88\% 以上 \cite{baselmsrefuse2024}，直接为我们所依赖的跨模型几何假设提供了实证基础。\textsc{Circuit Breakers} \cite{circuitbreakers2024} 通过 LoRA 适配器重定向有害表示。在 VLM 中，\textsc{Security Tensors} \cite{securitytensors2025} 识别出"安全层"的连续范围（LLaMA-3.2-11B-Vision 中为第 9--20 层），这些层的激活对有害文本输入响应，而对有害图像输入沉默；我们在 \textsc{OpenVLA} 的动作解码通路上复现了该探针。\textsc{CMRM} \cite{cmrm2025} 提供测试时位移向量；\textsc{VLM-Guard} \cite{vlmguard2025} 通过 SVD 从对齐 LLM 中提取 rank-$k$ 引导方向，并在推理时将其施加于 VLM 的最终输入 token 位置。\textsc{ASTRA} \cite{astra2025} 从视觉对抗样本中推导引导方向。\textsc{SafeSteer} \cite{safesteer2025} 将 SVD 子空间投影应用于 VLM 的通用越狱防御。

与我们工作最相近的是 \textsc{SARSteer} \cite{sarsteer2025}，该同期工作在音频语言模型中引导文本衍生的拒绝方向。关键区别在于：\textsc{SARSteer} 从目标模型自身提取方向，其主干已具备拒绝响应。而本文目标 VLA 建立在 \emph{base} LLM 之上，完全没有拒绝响应可供提取——在此场景下 \textsc{SARSteer} 的方法零样本不可用。我们从外部对齐的兄弟模型中提取方向并跨模型移植。

\paragraph{VLA 的激活引导用于行为控制。}
独立于安全问题，近期可解释性工作表明，激活引导对于 VLA 的\emph{行为}控制是有效的：\citet{mechinterpvla2025} 展示了在 \textsc{OpenVLA} 和 \textsc{$\pi_0$} 的特定 FFN 位置添加学习到的引导向量，可以可靠地修改运动速度和轨迹曲率。这为我们的安全干预建立了前提——\textsc{OpenVLA} 的残差流具有可被方向引导的几何结构——尽管其目标是行为控制而非危害拒绝，方向来源是分布内 VLA 数据而非跨模型拒绝方向，且未限制注入到动作 token 位置。

\paragraph{VLA 模型及其安全性。}
\textsc{OpenVLA} \cite{openvla2024} 和 \textsc{OpenVLA-OFT} \cite{openvlaoft2025} 代表了两种主流动作解码器：覆写 Llama 词表的离散动作 token，以及连续 L1 回归头。\textsc{$\pi_0$} \cite{pi02024} 在连续动作上使用流匹配。攻击性工作 \cite{wu2024exploringadversarial,badrobot2025,shawshank2025} 揭示了一个一致的结构性观察：VLA 没有拒绝训练的类比，且通用对抗补丁可跨任务迁移。防御方面，\textsc{SafeVLA} \cite{safevla2025} 将问题构建为约束马尔可夫决策过程，在安全强化学习下微调策略；\textsc{VLSA}/\textsc{AEGIS} \cite{vlsa2025} 在动作输出层施加控制论碰撞规避。这两种方法均未在表示层进行干预，即都没有询问 LLM 主干在生成动作之前是否本可感知该请求是有害的。

\paragraph{我们的定位。}
表~\ref{tab:positioning} 沿四个维度总结了设计空间：干预目标、阶段、方向来源和作用 token 位置。据我们所知，本工作是首个将安全方向从\emph{外部}对齐 LLM 提取并移植至以 \emph{base} 主干构建的 VLA 中，并限制注入到动作生成相关位置的方法，同时提供了覆盖完整 7 token 动作序列的生成阶段感知双阶段变体（\S\ref{sec:method:dualphase}）。

\begin{table*}[t]
\centering\small
\setlength{\tabcolsep}{5pt}
\begin{tabular}{p{4.4cm}p{2.0cm}p{0.8cm}p{2.8cm}p{2.4cm}}
\toprule
方法 & 目标 & 阶段 & 方向来源 & 作用位置 \\
\midrule
AdaShield \cite{adashield2024}        & 提示             & 推理 & ---               & N/A \\
TGA/CMRT-LVLM \cite{cmrt2025}         & 投影层 + LLM     & 训练 & 标题引导          & 全部 \\
CMRM \cite{cmrm2025}                  & LLM 隐层         & 推理 & 自身              & 全部 \\
VLM-Guard \cite{vlmguard2025}         & LLM 隐层         & 推理 & 外部 LLM          & 最后输入 token \\
ASTRA \cite{astra2025}                & LLM 隐层         & 推理 & 视觉对抗样本      & 生成 token \\
Security Tensors \cite{securitytensors2025} & 软 token    & 训练 & 有害 SFT          & 前缀 \\
SARSteer \cite{sarsteer2025}          & LLM 隐层         & 推理 & 自身              & 生成 token \\
SafeVLA \cite{safevla2025}            & 策略             & 训练 & 安全强化学习      & 全部 \\
Mech.\ Interp.\ VLA \cite{mechinterpvla2025} & LLM 隐层 & 推理 & 分布内 VLA 数据   & 全部 \\
\textbf{RDT/RDT+（本文）}             & LLM 隐层         & 推理 & \textbf{外部 LLM（跨模型）} & \textbf{动作位置 / 文本+动作} \\
\bottomrule
\end{tabular}
\caption{多模态安全设计空间中的定位。RDT/RDT+ 首次将（a）从\emph{外部}对齐 LLM 提取的方向移植至\emph{未对齐}的 base 主干 VLA，并（b）限制注入到动作生成位置。}
\label{tab:positioning_zh}
\end{table*}
